\begin{lemma}[The frame travel bound on a mark's portion]
Let $F$ be a finite-dimensional real inner-product space, let
$(\Omega,\mathcal M,\mu)$ be a measure space, let $(\mathcal G_n)_{n\in\N}$ be a
filtration of $\mathcal M$, let
$z,z^+,z^-:\N\times\Omega\to\ell^2(F\times\R)$, let $Q_0:\N\times\Omega\to F$,
let $S:\N\times\Omega\to\Z$, let $\sigma:\N\times\Omega\to\N$, let
$p:\N\times\Omega\to F$, let $\varsigma,\tau:\Omega\to\N$, let $j,b\in\N$, let
$\epsilon,\rho,C_p,B_0,C',C_M,C_2,L_K,\Lambda\in\R$ and let $u:\Z\to\R$. Write
$z_n=(x_n,t_n)$ and, for a stopping time $\varrho$, a slope field $q$ and a
scale $\rho$,
\[
  \mathcal T[\varrho,\varsigma,q]
   =\sup_{\varrho\le n\le\varsigma}
     \frac{\bigl|x_n-x_\varrho-(t_n-t_\varrho)\,q\bigr|}{\rho}
\]
for the running displacement in the sheared frame of slope $q$ and scale
$\rho$. Assume:
\begin{enumerate}
\item\label{it:mftb-pkg} the travel package of
\noderef{StoppedWalkTravel}: for all $(\mathcal G_n)$-stopping times
$\varrho\le\varsigma'\le\tau$, every $\rho'>0$, every $C\in\R$ and every
$q:\Omega\to F$ with $|q|\le C$ such that $\mathbf1_{\{\varrho\le n\}}q$ is
$\mathcal G_n$-measurable for all $n$, the four functions
\[
  \mathcal T[\varrho,\varsigma',q],\qquad
  \sum_{\varrho\le n<\varsigma'}\Bigl(\frac{\epsilon}{\rho'}\Bigr)^2|Q_0(n)-q|^2,
  \qquad
  \Bigl(\frac{\epsilon}{\rho'}\Bigr)^2(\varsigma'-\varrho),\qquad
  \sum_{\varrho\le n<\varsigma'}\mathcal E_n
\]
are $\mu$-integrable, where
$\mathcal E_n=|\epsilon^{-1}(x^+_n-x_n)-Q_0(n)|^2
+|\epsilon^{-1}(x_n-x^-_n)-Q_0(n)|^2$, and
\[
  \int\mathcal T[\varrho,\varsigma',q]\,d\mu
  \le2\Bigl(\int\sum_{\varrho\le n<\varsigma'}
      \tfrac{\epsilon^2}{\rho'^2}|Q_0(n)-q|^2d\mu\Bigr)^{1/2}
   +\Bigl(\int\tfrac{\epsilon^2}{\rho'^2}(\varsigma'-\varrho)\,d\mu\Bigr)^{1/2}
    \Bigl(\int\sum_{\varrho\le n<\varsigma'}\mathcal E_n\,d\mu\Bigr)^{1/2};
\]
\item\label{it:mftb-slope} $|p_j(\omega)|\le C_p$ for every $\omega$, and
$\mathbf1_{\{\sigma_j\le n\}}p_j$ is $\mathcal G_n$-measurable for every $n$;
\item\label{it:mftb-st} $\sigma_j$, $\varsigma$ and $\tau$ are
$(\mathcal G_n)$-stopping times with $\sigma_j\le\varsigma\le\tau$ pointwise;
\item\label{it:mftb-int} the portion length
$\omega\mapsto\varsigma(\omega)-\sigma_j(\omega)$ and the portion weight sum
$\omega\mapsto\sum_{\sigma_j\le m<\varsigma}u(S_m)$ are $\mu$-integrable, and
$\displaystyle\int\sum_{\sigma_j\le m<\varsigma}u(S_m)\,d\mu\le4\Lambda$;
\item\label{it:mftb-uw} $u(S_n(\omega))
=\bigl(\tfrac{\epsilon}{\rho}\bigr)^2
\bigl|\log\bigl(\tfrac{\epsilon}{\rho}S_n(\omega)\bigr)\bigr|^2$ whenever
$\sigma_j(\omega)\le n<\varsigma(\omega)$;
\item\label{it:mftb-Q} if $t_{\sigma_j(\omega)}(\omega)=\rho$ and
$\sigma_j(\omega)\le n<\varsigma(\omega)$ then
$|Q_0(n)(\omega)-p_j(\omega)|
\le C_2\bigl(L_K+\bigl|\log\bigl(\tfrac{\epsilon}{\rho}S_n(\omega)\bigr)\bigr|\bigr)$;
\item\label{it:mftb-off} if $t_{\sigma_j(\omega)}(\omega)\ne\rho$ then
$\varsigma(\omega)=\sigma_j(\omega)$;
\item\label{it:mftb-time} $\displaystyle
\int\bigl(\varsigma-\sigma_j\bigr)d\mu\le\frac{b^2}{4}$;
\item\label{it:mftb-budget} $\displaystyle
\int\sum_{n<\tau}\mathcal E_n\,d\mu\le B_0$;
\item\label{it:mftb-const} $0<\rho$, $\dfrac{\epsilon}{\rho}\,b=8$,
$C_M=2C_2^2(16+4\Lambda)$, $0\le C_M$, $0<C_2$, $0\le\Lambda$,
$2\sqrt{C_M}\le C'$, $4\le C'$ and $1\le L_K$.
\end{enumerate}
Then $\mathcal T[\sigma_j,\varsigma,p_j]$ is $\mu$-integrable and
\begin{equation}\label{eq:mftb-travel}
  \int\mathcal T[\sigma_j,\varsigma,p_j]\,d\mu
  \le C'\bigl(L_K+\sqrt{B_0}\bigr).
\end{equation}
\end{lemma}
\begin{proof}
\emph{The two instantiations.} Apply \ref{it:mftb-pkg} at
$\varrho=\sigma_j$, $\varsigma'=\varsigma$, $q=p_j$, $C=C_p$ and
$\rho'=\rho$; its hypotheses are \ref{it:mftb-st}, \ref{it:mftb-slope} and
$\rho>0$ from \ref{it:mftb-const}. This gives the four integrability facts for
the portion --- the first of them is the integrability asserted in the
conclusion --- and the travel inequality
\begin{equation}\label{eq:mftb-cs}
  \int\mathcal T[\sigma_j,\varsigma,p_j]\,d\mu
  \le2\sqrt{V}+\sqrt{T}\,\sqrt{E},
\end{equation}
where $V$, $T$ and $E$ are the three integrals on its right at this
instantiation. Apply \ref{it:mftb-pkg} a second time at
$\varrho=0$, $\varsigma'=\tau$, $q=0$, $C=0$ and $\rho'=1$ --- the zero map is
trivially bounded by $0$ and $\mathbf1_{\{0\le n\}}0$ is measurable --- and keep
only its fourth integrability fact, the integrability of
$\omega\mapsto\sum_{n<\tau}\mathcal E_n$.

\emph{The velocity term: $V\le C_ML_K^2$.} Fix $\omega$. If
$t_{\sigma_j(\omega)}(\omega)\ne\rho$ then $\varsigma(\omega)=\sigma_j(\omega)$
by \ref{it:mftb-off}, the summation range is empty and the summand bound below
holds trivially, both sides being $0$ on the left and nonnegative on the right.
Otherwise, for $\sigma_j(\omega)\le m<\varsigma(\omega)$, \ref{it:mftb-Q} gives
$|Q_0(m)-p_j|\le C_2(L_K+|\log(\tfrac\epsilon\rho S_m)|)$, so squaring and using
$(a+b)^2\le2a^2+2b^2$,
\[
  |Q_0(m)-p_j|^2
   \le C_2^2\Bigl(2L_K^2+2\bigl|\log\bigl(\tfrac\epsilon\rho S_m\bigr)\bigr|^2\Bigr),
\]
and multiplying by $(\epsilon/\rho)^2\ge0$ and using \ref{it:mftb-uw},
\[
  \Bigl(\frac\epsilon\rho\Bigr)^2|Q_0(m)-p_j|^2
   \le2C_2^2L_K^2\Bigl(\frac\epsilon\rho\Bigr)^2+2C_2^2\,u(S_m).
\]
Summing over $\sigma_j(\omega)\le m<\varsigma(\omega)$ and splitting the
constant term off,
\[
  \sum_{\sigma_j\le m<\varsigma}\Bigl(\frac\epsilon\rho\Bigr)^2|Q_0(m)-p_j|^2
   \le2C_2^2L_K^2\Bigl(\frac\epsilon\rho\Bigr)^2(\varsigma-\sigma_j)
     +2C_2^2\sum_{\sigma_j\le m<\varsigma}u(S_m).
\]
Both sides are integrable --- the left by the first instantiation, the right by
\ref{it:mftb-int} --- so integrating preserves the inequality. For the first
term, \ref{it:mftb-time} and $(\epsilon/\rho)^2\ge0$ give
$(\epsilon/\rho)^2\int(\varsigma-\sigma_j)\,d\mu\le(\epsilon/\rho)^2b^2/4$,
which is $16$ because $\bigl((\epsilon/\rho)b\bigr)^2=64$ by
\ref{it:mftb-const}; for the second, \ref{it:mftb-int} gives at most
$2C_2^2\cdot4\Lambda$. Hence, using $1\le L_K^2$ (from $1\le L_K$) to absorb the
second term into a multiple of $L_K^2$ and $0\le8C_2^2\Lambda$,
\[
  V\le2C_2^2L_K^2\cdot16+8C_2^2\Lambda
   \le2C_2^2(16+4\Lambda)L_K^2=C_ML_K^2
\]
by the definition of $C_M$ in \ref{it:mftb-const}.

\emph{The error term: $E\le B_0$.} Pointwise, the portion's index range
$[\sigma_j,\varsigma)$ is contained in $[0,\tau)$ by \ref{it:mftb-st}, and each
summand $\mathcal E_n$ is nonnegative, so
$\sum_{\sigma_j\le n<\varsigma}\mathcal E_n\le\sum_{n<\tau}\mathcal E_n$. Both
sides are integrable --- the left by the first instantiation, the right by the
second --- so $E\le\int\sum_{n<\tau}\mathcal E_n\,d\mu\le B_0$ by
\ref{it:mftb-budget}.

\emph{The time term: $T\le16$.} Since $\sigma_j\le\varsigma$ pointwise, the
real-valued difference $\varsigma-\sigma_j$ agrees with the cast of the natural
subtraction, so
$T=(\epsilon/\rho)^2\int(\varsigma-\sigma_j)\,d\mu\le(\epsilon/\rho)^2b^2/4=16$
exactly as in the velocity term.

\emph{Assembly.} The square root is monotone, so
$\sqrt V\le\sqrt{C_ML_K^2}=\sqrt{C_M}\,L_K$ --- using $L_K\ge1>0$ ---
$\sqrt T\le\sqrt{16}=4$ and $\sqrt E\le\sqrt{B_0}$. Since
$\sqrt T\ge0$ and $\sqrt E\ge0$, the product is bounded by
$4\sqrt{B_0}$. Substituting into \eqref{eq:mftb-cs},
\[
  \int\mathcal T[\sigma_j,\varsigma,p_j]\,d\mu
   \le2\sqrt{C_M}\,L_K+4\sqrt{B_0}
   \le C'L_K+C'\sqrt{B_0}
   =C'\bigl(L_K+\sqrt{B_0}\bigr),
\]
using $2\sqrt{C_M}\le C'$ with $L_K\ge0$ and $4\le C'$ with
$\sqrt{B_0}\ge0$. This is \eqref{eq:mftb-travel}.
\end{proof}
